	\documentclass[12pt]{exam}
	\usepackage[utf8]{inputenc}
	\usepackage[portuguese,brazil]{babel}
	\usepackage[margin=.65in]{geometry}
	\usepackage{amsmath,amssymb}
	\usepackage{multicol}
	\usepackage{multirow}
	\usepackage{float}
	\usepackage{enumerate}
	\usepackage{enumitem}
	\usepackage{array}
	\usepackage{graphicx}
	\usepackage{listings}
	\usepackage{circuitikz}    
	\usepackage{xcolor}
	\usepackage{tikz}
	\usetikzlibrary{shapes.multipart, calc}
	\usepackage{rotating}
	\usepackage{caption}
	%\usepackage[version=3]{mhchem}

	\pointpoints{ponto}{pontos}
	\hpword{Pontos:}
	\vpword{Pontos}
	\htword{Total}
	\vtword{Total:}
	\vsword{Resultado}

	\definecolor{azul}{RGB}{0,32,96}
	\definecolor{verde}{RGB}{0,132,124}
	\definecolor{cinza}{RGB}{0,84,132}
	%\definecolor{cinza2}{RGB}{218,220,216}
	\definecolor{cinza2}{RGB}{132,134,136}
	\definecolor{mygreen}{rgb}{0,0.6,0}
	\definecolor{mygray}{rgb}{0.5,0.5,0.5}
	\definecolor{mymauve}{rgb}{0.58,0,0.82}
	\newcommand*{\vttfamily}{%
	\fontencoding{T1}\fontfamily{lmvtt}\selectfont
	}
	\lstdefinestyle{Cstyle}{%
	language=C,numbers=left, keywordstyle=\color{mymauve}\vttfamily,commentstyle=\color{mygreen}\vttfamily,stringstyle=\color{blue},basicstyle=\scriptsize\vttfamily,numberstyle=\tiny\vttfamily,breaklines=true,stepnumber=1,tabsize=2,xleftmargin=\leftmarginii,escapechar=!,captionpos=b,morekeywords={include, printf,NULL},columns=flexible,escapeinside={(^}{^)},
	frame=l,
	framesep=2.5mm,
	framexleftmargin=2.5mm,
	fillcolor=\color{cinza2!50},
	rulecolor=\color{ballblue},
	rulecolor=\color{blue}
	}

	\lstnewenvironment{C}[1][]
	{
	\lstset{style=Cstyle,#1}
	}
	{}


	\newcommand{\class}{Segurança da Informação}
	\newcommand{\profname}{Gilmar Gomes do Nascimento}
	\newcommand{\examnum}{Exercício - sessão, IDOR} \newcommand{\examdate}{$\today$}
	\newcommand*{\oldneg}{\mathord{\nao}}
	\pagestyle{headandfoot}
	\runningheadrule
	\firstpageheader{}{}{}
	\runningheader{\class}{\examnum}{\examdate}
	\runningfootrule
	\firstpagefooter{}{}{}
	\runningfooter{}{Página \thepage\ de \numpages}{}
	%\pagestyle{foot}
	%\lfoot{}
	%\cfoot[]{Página \thepage\ de \numpages}
	%\rfoot{}
	%
	%\runningfootrule
	%\boxedpoints %\pointsinmargin
	%\printanswers
	%\SolutionEmphasis{\color{red}}
	%\definecolor{SolutionBoxColor}{gray}{0.8}
	%\renewcommand{\solutiontitle}{\noindent\textbf{Solução:}\enspace}
	%\renewcommand{\solutiontitle}{\noindent\textbf{Solution:}\par\noindent}

	\begin{document}

	\noindent
	\begin{tabular*}{\textwidth}{l @{\extracolsep{\fill}} r @{\extracolsep{6pt}} l}
	\textbf{Disciplina:} \class &  \multicolumn{2}{c}{\multirow{4}{*}{
	\includegraphics[scale=0.35]{logo.png}}}\\
	\textbf{Prof.:} ~\profname \vspace{0.2in} && \\
	\vspace{0.2in} \textbf{Nome:} \makebox[3in]{\hrulefill} \vspace{0.2in} &&\\
	%\textbf{Matrícula:} \makebox[1.5in]{\hrulefill} \textbf{Nota:} \makebox[1in]{\hrulefill} &&\\ 
	%\multicolumn{3}{c}{\textbf{{\examnum \hspace{3px} \examdate}}}\\
	\multicolumn{3}{c}{}\\
	\end{tabular*}\\
	\rule[2ex]{\textwidth}{1.6pt}
	\parskip=0.1in

	%Este exame contêm \numpages\ páginas (incluindo esta capa) e \numquestions\ questões.\\
	%Total de \numpoints~ pontos.

	%\begin{center}
	%Grade Table (for teacher use only)\\
	%\addpoints
	%\gradetable[v][questions]
	%\end{center}

	\noindent

\begin{questions}

\question O que é uma \textbf{sessão} em aplicações web? Explique o papel do cookie de sessão e por que ele é usado em vez de enviar usuário e senha em cada requisição.
\begin{solutionbox}{2in}
\end{solutionbox}

\question O que é \textbf{entropia de token}? Por que ela é fundamental para a segurança de tokens de sessão? O que acontece se um token tiver baixa entropia?
\begin{solutionbox}{1.7in}
\end{solutionbox}

\question Explique o que é a flag \texttt{HttpOnly} em cookies. Qual ataque ela previne e por quê?
\begin{solutionbox}{2in}
\end{solutionbox}

\question Explique o que é a flag \texttt{Secure} em cookies. Em que cenário a ausência dessa flag representa risco?
\begin{solutionbox}{1.7in}
\end{solutionbox}

\question Uma aplicação gera tokens de sessão usando apenas \texttt{timestamp} em segundos. Explique por que isso é inseguro e como um atacante poderia explorar essa falha.
\begin{solutionbox}{1.7in}
\end{solutionbox}

% ===== BLOCO 2: PRINCÍPIOS DE SEGURANÇA =====

%\section*{Bloco 2: Princípios de Segurança}

\question O que significa o princípio \textbf{Never Trust the Client} (nunca confie no cliente)? Dê um exemplo prático de falha causada pela violação desse princípio.
\begin{solutionbox}{1.7in}
\end{solutionbox}

\question Explique o conceito de \textbf{Source of Truth} (fonte da verdade) no contexto de segurança de aplicações. Por que o servidor deve ser sempre a fonte da verdade?
\begin{solutionbox}{1.7in}
\end{solutionbox}

%\question Relacione o princípio \textbf{Never Trust the Client} com o lab de IDOR que você resolveu. Como a falha poderia ter sido evitada se o servidor aplicasse esse princípio corretamente?
%\begin{solutionbox}{3in}
%\end{solutionbox}

%\question Analise a situação: uma aplicação web valida o preço de um produto no JavaScript do navegador antes de enviar ao servidor. Explique por que isso é uma falha de segurança e qual seria a abordagem correta.
%\begin{solutionbox}{2.5in}
%\end{solutionbox}

% ===== BLOCO 3: FUZZING E AUTOMAÇÃO =====

%\section*{Bloco 3: Fuzzing e Automação com Burp}

\question O que é \textbf{fuzzing} no contexto de testes de segurança? Qual o objetivo dessa técnica?
\begin{solutionbox}{2.5in}
\end{solutionbox}

\question Explique a diferença entre os três tipos de ataque do Burp Intruder:
\begin{itemize}
\item \textbf{Sniper}
\item \textbf{Pitchfork}
\item \textbf{Cluster Bomb}
\end{itemize}
\begin{solutionbox}{3.2in}
\end{solutionbox}

%\question Em que cenário você usaria o modo \textbf{Sniper}? E em que cenário o \textbf{Cluster Bomb} seria mais adequado que o \textbf{Pitchfork}?
%\begin{solutionbox}{3in}
%\end{solutionbox}
%
%\question Considere que você precisa testar um formulário de login com 100 usuários e 100 senhas. Qual modo do Intruder você usaria para testar todas as combinações possíveis? Justifique.
%\begin{solutionbox}{2.5in}
%\end{solutionbox}
%
%\question Explique como o fuzzing pode ser usado para descobrir \textbf{tokens de sessão previsíveis} ou \textbf{IDs de recursos enumeráveis}.
%\begin{solutionbox}{3in}
%\end{solutionbox}

\end{questions}

\end{document}
